\section{Subfamily theorems}\label{sec:subfamily}

This section proves $\Lprime$ unconditionally on books. On $2$-paths
we prove the Ces\`aro square-energy identity
$\sminus(\Ln)/n \to c_{1}^{-}$ (\Cref{thm:2path-szego}); the
companion ear-gain statement $\dminus(\Ln) \to c_{1}^{-}$ is
\emph{not} proved at the asymptotic first-difference level (open
subobligation O13.1, \Cref{rem:2path-first-difference}). For
$n \in [4, 1000]$ we give an engineering-rigorous finite-range
floor $\dminus(\Ln) \ge 21/16$ via a Demmel--Kahan a-posteriori
certificate, conditional on a conservative \texttt{dsyevd}
backward-error envelope (\Cref{thm:2path-DK}). On the $\BTkt{k}{2}$
family we exhibit the bad-tail-ear asymptotic that refutes the
\emph{universal} ear-deletion form of $\Lprime$
(\Cref{thm:BT-asymptotic}, \Cref{cor:universal-false}). The
existential $\Lprime$ on $\BTkt{k}{2}$ (via the max-degsum selector,
which picks a clean book-page ear) is verified empirically on the
working corpus; we do not claim an analytical closure of $\Lprime$ on
the full $\BTkt{k}{2}$ family in this paper. The Stieltjes-transform
identification of the $2$-path negative-spectrum measure is in
\Cref{sec:subfamily:2path-stieltjes}; its use in proving condition~(a)
of the joint-invariant ansatz on the $2$-path family is deferred to
\Cref{sec:moment-form:stieltjes}. (\emph{Caveat:} the $2$-path is
\emph{not} the universal binding case for $\Iv$ at finite~$n$; see F10
in \Cref{sec:failure-modes}.)

\subsection{Books \texorpdfstring{$\Bk$}{B\_k}}\label{sec:subfamily:books}

The book $\Bk$ (also written $K_{1, 1, k}$) is the $2$-tree on $k + 2$
vertices consisting of $k$ triangles sharing a common edge $\{a, b\}$;
its non-spine vertices $p_{1}, \ldots, p_{k}$ (the \emph{pages}) are
pairwise non-adjacent simplicial vertices of degree $2$, with supporting
edge $\{a, b\}$.

\begin{theorem}[Closed form on books]\label{thm:books}
For every $k \ge 2$, every page ear $v$ of $\Bk$ satisfies
\[
\dminus(\Bk) \;=\; 2 - \frac{4}{\sqrt{8k+1} + \sqrt{8k-7}},
\qquad
\dplus(\Bk) \;=\; 2 + \frac{4}{\sqrt{8k+1} + \sqrt{8k-7}}.
\]
Both gains lie in $[17/16, 47/16]$ for every $k \ge 2$. The minimum is
$\dminus(B_{2}) = (7 - \sqrt{17})/2 \approx 1.4385$ at $k = 2$, and
$\dminus(\Bk) \uparrow 2$ as $k \to \infty$.
\end{theorem}

\begin{proof}
We first compute the spectrum of $\Adj(\Bk)$. The book has the natural
symmetries $a \leftrightarrow b$ (spine swap) and $S_{k}$
(page permutation), inducing an orthogonal decomposition of
$\mathbb{R}^{k+2}$ into three invariant subspaces:
\begin{enumerate}
\item The \emph{page-antisymmetric} subspace
$\{x : x_{a} = x_{b} = 0, \sum_{i} x_{p_{i}} = 0\}$, of dimension
$k - 1$. For any such $x$, $(\Adj(\Bk) x)_{a} = x_{b} + \sum_{i} x_{p_{i}} = 0$,
$(\Adj(\Bk) x)_{b} = x_{a} + \sum_{i} x_{p_{i}} = 0$, and
$(\Adj(\Bk) x)_{p_{i}} = x_{a} + x_{b} = 0$. Hence this subspace lies
in $\ker \Adj(\Bk)$, contributing eigenvalue $0$ with multiplicity
$k - 1$.

\item The \emph{spine-antisymmetric} subspace
$\{x : x_{a} = -x_{b}, x_{p_{i}} = 0 \text{ for all } i\}$, of
dimension $1$. For $x = e_{a} - e_{b}$ we compute
$(\Adj(\Bk) x)_{a} = -1 = -x_{a}$, $(\Adj(\Bk) x)_{b} = 1 = -x_{b}$,
$(\Adj(\Bk) x)_{p_{i}} = 0$. So $\Adj(\Bk)(e_{a} - e_{b}) = -(e_{a} - e_{b})$,
giving eigenvalue $-1$ with multiplicity $1$.

\item The \emph{symmetric} subspace
$\{x : x_{a} = x_{b}, x_{p_{1}} = \cdots = x_{p_{k}}\}$, of
dimension $2$, spanned by $e_{a} + e_{b}$ and $\sum_{i} e_{p_{i}}$.
In the orthonormal basis
$u_{1} := (e_{a} + e_{b})/\sqrt{2}$,
$u_{2} := (\sum_{i} e_{p_{i}})/\sqrt{k}$, the restriction of
$\Adj(\Bk)$ is
$
\begin{pmatrix}
1 & \sqrt{2k} \\
\sqrt{2k} & 0
\end{pmatrix},
$
with characteristic polynomial $\lambda^{2} - \lambda - 2k = 0$ and
roots $\lambda_{\pm} = (1 \pm \sqrt{8k + 1})/2$.
\end{enumerate}
Combining,
\[
\spec(\Adj(\Bk)) \;=\;
\Bigl\{ \tfrac{1 + \sqrt{8k+1}}{2},\; \underbrace{0, \ldots, 0}_{k-1},\;
-1,\; \tfrac{1 - \sqrt{8k+1}}{2} \Bigr\}.
\]
For $k \ge 1$ both $-1$ and $(1 - \sqrt{8k+1})/2 \le -1$ are
negative, so
\[
\sminus(\Bk) \;=\; 1 + \left(\tfrac{1 - \sqrt{8k+1}}{2}\right)^{2}
\;=\; 1 + \tfrac{1}{4}\bigl(1 - 2\sqrt{8k+1} + 8k + 1\bigr)
\;=\; \tfrac{1}{2}\bigl(4k + 3 - \sqrt{8k+1}\bigr).
\]

The page-ear deletion $\Bk - v = B_{k-1}$ for any page vertex $v$, so
$\sminus(B_{k} - v) = \sminus(B_{k-1}) = \tfrac{1}{2}(4(k-1) + 3 - \sqrt{8(k-1)+1})
= \tfrac{1}{2}(4k - 1 - \sqrt{8k-7})$. Subtracting,
\[
\dminus(\Bk) \;=\; \sminus(\Bk) - \sminus(B_{k-1})
\;=\; \tfrac{1}{2}\bigl[(4k + 3) - (4k - 1) - \sqrt{8k+1} + \sqrt{8k-7}\bigr]
\;=\; 2 - \tfrac{1}{2}\bigl(\sqrt{8k+1} - \sqrt{8k-7}\bigr).
\]
Rationalising the difference of square roots,
$\sqrt{8k+1} - \sqrt{8k-7} = 8/(\sqrt{8k+1} + \sqrt{8k-7})$, giving
\[
\dminus(\Bk) \;=\; 2 - \frac{4}{\sqrt{8k+1} + \sqrt{8k-7}},
\]
as stated. The companion formula
$\dplus(\Bk) = 4 - \dminus(\Bk) = 2 + 4/(\sqrt{8k+1} + \sqrt{8k-7})$
follows from the trace identity $\dplus + \dminus = 2 \deg_{\Bk}(v) = 4$.

For monotonicity, $\sqrt{8k+1} + \sqrt{8k-7}$ is strictly increasing in
$k$, so $\dminus(\Bk)$ is strictly increasing and $\dminus(\Bk) \uparrow 2$
as $k \to \infty$. The minimum at $k = 2$ is
$\dminus(B_{2}) = 2 - 4/(\sqrt{17} + 3) = 2 - (4(\sqrt{17} - 3))/(17 - 9)
= 2 - (\sqrt{17} - 3)/2 = (7 - \sqrt{17})/2 \approx 1.4385$.

For the gain bounds in $[17/16, 47/16]$: $\dminus(\Bk)$ is bounded below
by its minimum $(7 - \sqrt{17})/2 \approx 1.4385 > 17/16 = 1.0625$ and
above by $\sup_{k} \dminus(\Bk) = 2 < 47/16 = 2.9375$. By the trace
identity, $\dplus(\Bk) = 4 - \dminus(\Bk)$ is bounded below by
$4 - 2 = 2 > 17/16$ (since $\dminus(\Bk) < 2$) and above by
$4 - \dminus(B_{2}) = (1 + \sqrt{17})/2 \approx 2.5616 < 47/16$. Hence
both gains lie in $[17/16, 47/16]$ for every $k \ge 2$.
\end{proof}

\begin{corollary}\label{cor:books-Lprime}
$\Lprime$ holds unconditionally on the book family $\Bk$ for every
$k \ge 2$, with any page vertex as the witness ear.
\end{corollary}

\begin{proof}
Every page vertex $v$ is a simplicial degree-$2$ ear of $\Bk$ and
satisfies $\dminus(\Bk) \ge 17/16$ and $\dplus(\Bk) \ge 17/16$ by
\Cref{thm:books}.
\end{proof}

\subsection{$2$-paths \texorpdfstring{$\Ln$}{L\_n}: the Szeg\H{o} limit}\label{sec:subfamily:2path-szego}

The $2$-path $\Ln$ is the square of the path graph $P_{n}$: the vertex
set is $\{1, \ldots, n\}$, and $i \sim j$ iff $|i - j| \in \{1, 2\}$.
The adjacency matrix $\Adj(\Ln)$ is symmetric pentadiagonal Toeplitz
with first row $(0, 1, 1, 0, \ldots, 0)$ and symbol
\[
f(\theta) \;=\; 2 \cos\theta + 2 \cos 2\theta \;=\; 2(2\cos\theta - 1)(\cos\theta + 1),
\]
the second equality from $\cos 2\theta = 2\cos^{2}\theta - 1$. The
zeros of $f$ on $[0, \pi]$ are $\theta = \pi/3$ and $\theta = \pi$, and
$f$ is negative on $(\pi/3, \pi)$. The maximum is $f(0) = 4$ and the
minimum is $f(\arccos(-1/4)) = -9/4$.

The two simplicial degree-$2$ ears of $\Ln$ are the endpoints $v = 1$
and $v = n$, which are mapped to each other by the reflection
$i \mapsto n + 1 - i$ and are therefore isomorphic. We write
$\dminus(\Ln) := \dminus(1)$ and $\dplus(\Ln) := \dplus(1)$.

\begin{theorem}[Szeg\H{o} Ces\`aro limit on $2$-paths]\label{thm:2path-szego}
With the notation above,
\[
\lim_{n \to \infty} \frac{\sminus(\Ln)}{n}
\;=\; \frac{32 \pi - 27 \sqrt{3}}{12 \pi}
\;=:\; c_{1}^{-}
\;\approx\; 1.4262,
\qquad
\lim_{n \to \infty} \frac{\splus(\Ln)}{n}
\;=\; \frac{16 \pi + 27 \sqrt{3}}{12 \pi}
\;=:\; c_{1}^{+}
\;\approx\; 2.5738,
\]
both strictly inside $(17/16, 47/16)$. Numerically the
first differences $\dminus(\Ln)$ converge to $c_{1}^{-}$ along the
sequence, attaining their (empirically verified) minimum $\approx 1.319$
at $n = 6$ before increasing monotonically toward $c_{1}^{-}$ (see
\Cref{rem:2path-first-difference}).
\end{theorem}

\begin{proof}
By Szeg\H{o}'s first theorem for banded Hermitian Toeplitz matrices
\cite{GrenanderSzego1958, BottcherSilbermann}, for every continuous
$\phi : \mathbb{R} \to \mathbb{R}$,
\[
\lim_{n \to \infty} \frac{1}{n} \tr\,\phi\bigl(\Adj(\Ln)\bigr)
\;=\; \frac{1}{2\pi}\int_{-\pi}^{\pi} \phi\bigl(f(\theta)\bigr)\, d\theta.
\]
The function $\phi(\lambda) = \lambda^{2} \mathbf{1}_{\lambda < 0}$ is
continuous on the bounded spectrum $[f_{\min}, f_{\max}] = [-9/4, 4]$
(the would-be discontinuity at $\lambda = 0$ is removable because
$\phi(0^{-}) = 0 = \phi(0) = \phi(0^{+})$). Szeg\H{o}'s first theorem
therefore gives
\[
\lim_{n \to \infty} \frac{\sminus(\Ln)}{n}
\;=\; \frac{1}{2\pi} \int_{-\pi}^{\pi} f(\theta)^{2}\,
\mathbf{1}_{f(\theta) < 0}\, d\theta
\;=\; \frac{1}{\pi} \int_{\pi/3}^{\pi} f(\theta)^{2}\, d\theta,
\]
where the second equality uses the evenness of $f$ and the sign
pattern $\mathbf{1}_{f < 0} = \mathbf{1}_{(\pi/3, \pi)}$ on $(0, \pi)$.

The theorem statement uses Ces\`aro limits, which are immediate from
the above. The first-difference convergence
$\dminus(\Ln) \to c_{1}^{-}$ does \emph{not} follow from the
leading-order Szeg\H{o} theorem alone (a sequence $a_{n}$ with
$a_{n}/n \to c$ need not satisfy $a_{n} - a_{n-1} \to c$); it is
treated separately in \Cref{rem:2path-first-difference} below.

Expand $f(\theta)^{2} = (2\cos\theta + 2\cos 2\theta)^{2}
= 4\cos^{2}\theta + 8\cos\theta\cos 2\theta + 4\cos^{2}2\theta$.
Using the product-to-sum identity
$\cos\theta \cos 2\theta = (\cos\theta + \cos 3\theta)/2$ and the
half-angle identities $2\cos^{2}\theta = 1 + \cos 2\theta$,
$2\cos^{2} 2\theta = 1 + \cos 4\theta$, we obtain
\[
f(\theta)^{2} \;=\; 4 + 2\cos 2\theta + 2\cos 4\theta + 4\cos\theta + 4\cos 3\theta.
\]
Term-by-term integration over $(\pi/3, \pi)$:
\[
\int_{\pi/3}^{\pi} 4\, d\theta = \tfrac{8\pi}{3}, \quad
\int_{\pi/3}^{\pi} 2\cos 2\theta\, d\theta = -\tfrac{\sqrt{3}}{2}, \quad
\int_{\pi/3}^{\pi} 2\cos 4\theta\, d\theta = \tfrac{\sqrt{3}}{4},
\]
\[
\int_{\pi/3}^{\pi} 4\cos\theta\, d\theta = -2\sqrt{3}, \quad
\int_{\pi/3}^{\pi} 4\cos 3\theta\, d\theta = 0.
\]
Summing,
\[
\int_{\pi/3}^{\pi} f(\theta)^{2}\, d\theta
\;=\; \tfrac{8\pi}{3} - \tfrac{\sqrt{3}}{2} + \tfrac{\sqrt{3}}{4} - 2\sqrt{3}
\;=\; \tfrac{8\pi}{3} - \tfrac{9\sqrt{3}}{4}
\;=\; \tfrac{32\pi - 27\sqrt{3}}{12}.
\]
Dividing by $\pi$ gives $c_{1}^{-} = (32\pi - 27\sqrt{3})/(12\pi)$.
The companion formula for $c_{1}^{+}$ follows from the per-step trace
identity $\dplus(\Ln) + \dminus(\Ln) = 2\deg_{\Ln}(1) = 4$ at every
finite $n \ge 2$, which Ces\`aro-averages to
$\splus(\Ln)/n + \sminus(\Ln)/n \to 4$ (a consistency check: this
also follows from the Szeg\H{o}--Parseval identity, since
$f(\theta) = e^{i\theta} + e^{-i\theta} + e^{2i\theta} + e^{-2i\theta}$
has exactly four unit Fourier coefficients, so
$(1/\pi) \int_{0}^{\pi} f(\theta)^{2}\, d\theta = \sum |\hat{f}(k)|^{2} = 4$).
Hence $c_{1}^{+} = 4 - c_{1}^{-} = (16\pi + 27\sqrt{3})/(12\pi)$.

Both Ces\`aro limits lie strictly inside $(17/16, 47/16) \approx (1.063, 2.938)$.
\end{proof}

\begin{remark}[First-difference convergence $\dminus(\Ln) \to c_{1}^{-}$]
\label{rem:2path-first-difference}
The first-difference convergence $\dminus(\Ln) \to c_{1}^{-}$ is the
sharper companion statement to the Ces\`aro identity proved above. It
is verified in this paper at three independent levels.
\begin{enumerate}[label=(\roman*)]
\item \emph{Empirically.} The harness
\texttt{tests/test\_two\_path\_widom\_tightness.py} computes
$\dminus(\Ln)$ at every $n \in [4, 200]$ via direct \texttt{eigvalsh}
and asserts $|\dminus(L_{200}) - c_{1}^{-}| < 10^{-3}$, with the
sequence approaching $c_{1}^{-}$ monotonically from below after its
worst case $\dminus(L_{6}) \approx 1.3190$.
\item \emph{Finitely-rigorously up to $n = 1000$.} The
Demmel--Kahan certificate of \Cref{thm:2path-DK} gives the bound
$\dminus(\Ln) \ge 17/16 + 1/4 = 21/16$ uniformly across
$n \in [4, 1000]$. \emph{Warning on the meaning of ``slack $0.257$''.}
The number $0.257$ that appears in the proof of \Cref{thm:2path-DK}
is the gap between the empirical worst case
$\dminus(L_{6}) \approx 1.3187$ and $17/16 = 1.0625$
(i.e.\ $1.3187 - 17/16 \approx 0.2565$). It is the
slack of the \emph{empirical} $\dminus$ above $17/16$, used in the
proof solely to dominate the propagated forward-error bound
$1.42 \cdot 10^{-4}$ by three orders of magnitude --- which is what
licenses concluding $\dminus(\Ln) > 17/16 + 1/4 = 21/16$.
It is \emph{not} a safety margin above $21/16$: the gap
$\dminus(L_{6}) - 21/16 \approx 0.0062$ is two orders of magnitude
smaller, and the declared ``$+1/4$'' is the target offset of the
certificate, not an empirical buffer above $21/16$. The propagated
forward-error bound on $\dminus$ is $\le 1.42 \cdot 10^{-4}$.
\item \emph{Asymptotically.} The bordering identity
$\Adj(\Ln) = \bigl(\begin{smallmatrix} 0 & \earw^{\top} \\ \earw &
\Adj(L_{n-1}) \end{smallmatrix}\bigr)$ with $\earw = e_{1} + e_{2}$,
together with the weak convergence
$\mu_{w, n} \to \mu_{w, \infty}$ proved in
\Cref{thm:2path-stieltjes} and the no-atom-at-$0$ property of
$\mu_{w, \infty}$ (\Cref{thm:2path-stieltjes}, Step~4), heuristically
identifies the limit as the second-moment integral of
$\mu_{w, \infty}$ on the negative branch. A fully rigorous limit-value
proof at the first-difference level requires either a
Bogoya--B\"ottcher--Grudsky-type finite-$n$ expansion extended past
the simple-loop hypothesis (F9, \Cref{rem:simple-loop}) or a
Birman--Yafaev / Krein--Lifshitz spectral-shift-function argument
\cite{BirmanYafaev1993} for the rank-one bordering; neither is
necessary for any later argument in this paper, since the closure of
$\Lprime$ on $2$-paths at finite~$n$ is delivered by
\Cref{thm:2path-DK} (Demmel--Kahan) and the closure of condition~(a)
in the asymptotic limit is delivered by \Cref{thm:moment-form-stieltjes}
via the spectral measure $\mu_{w, \infty}$, not via $\dminus(\Ln)$.
We therefore record the first-difference convergence
$\dminus(\Ln) \to c_{1}^{-}$ as a proposition-level statement supported
by (i)--(ii) above with the asymptotic identification at
the level of $\sminus(\Ln)/n$; the formal first-difference limit proof
is listed as open subobligation O13.1 of \Cref{rem:simple-loop}.
\end{enumerate}
\end{remark}

\begin{remark}[Failure mode F9: simple-loop hypothesis]\label{rem:simple-loop}
The standard Bogoya--B\"ottcher--Grudsky \cite{BBG2015} /
Avram--Parter \cite{AvramParter1988} / Widom-type effective constants
for finite-$n$ corrections to Szeg\H{o}'s theorem require a
\emph{simple-loop} symbol — $f$ traces the unit circle once without
singularities. The pentadiagonal symbol
$f(\theta) = 2\cos\theta + 2\cos 2\theta$ has zeros at \emph{both} the
interior point $\theta = \pi/3$ and the boundary $\theta = \pi$, so
the simple-loop hypothesis fails. This affects the convergence
\emph{rate} in \Cref{thm:2path-szego} (recorded as open subobligation
O13.1 in \Cref{sec:open-problems}) but does not affect the limit value
displayed above, which follows from the unmodified Szeg\H{o} theorem.
\end{remark}

\subsection{$\BTkt{k}{2}$: a structured counterexample}\label{sec:subfamily:BT}

The family $\BTkt{k}{2}$ is the $2$-tree on $n = k + 4$ vertices with:
\begin{itemize}
\item book-spine vertices $0, 1$,
\item book pages $2, 3, \ldots, k+1$ (each adjacent to both $0$ and $1$),
\item first tail vertex $k+2$, adjacent to $0$ and $2$ (forming the
triangle $\{0, 2, k+2\}$ on the supporting edge $\{0, 2\}$),
\item second tail vertex $k+3$, adjacent to $2$ and $k+2$ (forming the
triangle $\{2, k+2, k+3\}$ on the supporting edge $\{2, k+2\}$).
\end{itemize}
The clique tree of $\BTkt{k}{2}$ consists of a star of $k$ ``page''
triangles $\{0, 1, j\}$ for $j \in \{2, \ldots, k+1\}$ around the
spine edge $\{0, 1\}$, plus a length-$2$ tail chain
$\{0, 2, k+2\}, \{2, k+2, k+3\}$ glued at the page triangle
$\{0, 1, 2\}$ through the shared edge $\{0, 2\}$. The tail simplicial
ear is $v_{\mathrm{tail}} := k + 3$, with neighbours $\{2, k+2\}$.

\begin{proposition}[Bad-ear asymptotic on $\BTkt{k}{2}$ -- computational / asymptotic evidence; conditional on the secular-rate subobligation]
\label{thm:BT-asymptotic}
With the notation above, conditional on the secular-equation rate
subobligation $a_{\pm} - b_{\pm} = O(1/\sqrt{k})$ (the
leading-eigenvalue square-cancellation step that controls the
$\sqrt{2k}$ contribution from the embedded book block; sketched in the
proof below),
\[
\lim_{k \to \infty} \dminus(v_{\mathrm{tail}})
\;=\; 4 - \alpha^{2} + \beta^{2}
\;\approx\; 1.0353,
\]
where $\alpha$ is the unique \emph{positive} root of
$P_{G}(x) := 2x^{3} - 7x - 3 = 0$ ($\alpha \approx 2.0565$) and
$\beta$ is the unique \emph{positive} root of
$P_{H}(x) := 2x^{3} + 2x^{2} - 3x - 2 = 0$ ($\beta \approx 1.1246$).
Each cubic has three distinct real roots; we take the largest in each
case. The equitable-partition reduction to the $6 \times 6$ (resp.\
$5 \times 5$) quotient matrix, the cubic interior limits, and the
numerical limit value are derived in full; the leading-eigenvalue
square-cancellation step is presented at the level of a secular
equation but not fully closed (see the proof below). The full
derivation is regression-locked numerically by
\texttt{tests/test\_lprime\_subfamilies.py}.
\end{proposition}

\begin{proof}[Proof sketch]
The non-tail pages $3, 4, \ldots, k+1$ of $\BTkt{k}{2}$ are pairwise
interchangeable by an $S_{k-1}$ symmetry (each is adjacent only to
$\{0, 1\}$, with no triangle through any pair). Set $m := k - 1$ and
decompose $\mathbb{R}^{k+4}$ into the $(m - 1)$-dimensional
\emph{page-trace-zero} subspace
$\{x : x_{0} = x_{1} = x_{2} = x_{k+2} = x_{k+3} = 0,\ \sum_{j=3}^{k+1} x_{j} = 0\}$,
which lies in $\ker \Adj(\BTkt{k}{2})$ by the same argument as in
\Cref{thm:books}, and its $6$-dimensional orthogonal complement
spanned by
\[
e_{0},\quad e_{1},\quad e_{2},\quad
u_{P} := m^{-1/2} \sum_{j=3}^{k+1} e_{j},\quad
e_{k+2},\quad e_{k+3}.
\]
In this basis $\Adj(\BTkt{k}{2})$ restricts to the symmetric
$6 \times 6$ matrix
\[
M_{G}(m)
\;=\;
\begin{pmatrix}
0 & 1 & 1 & \sqrt{m} & 1 & 0 \\
1 & 0 & 1 & \sqrt{m} & 0 & 0 \\
1 & 1 & 0 & 0 & 1 & 1 \\
\sqrt{m} & \sqrt{m} & 0 & 0 & 0 & 0 \\
1 & 0 & 1 & 0 & 0 & 1 \\
0 & 0 & 1 & 0 & 1 & 0
\end{pmatrix},
\]
encoding the spine edge $\{0, 1\}$ (entry $M_{01} = 1$), the page
adjacencies of vertex $2$ to both spines (entries $M_{02} = M_{12} = 1$),
the page-cluster couplings $\sqrt{m}$ (entries $M_{03} = M_{13} = \sqrt{m}$
from $\Adj e_{0} \cdot u_{P} = m^{-1/2} \sum_{j=3}^{k+1} 1 = \sqrt{m}$),
the first tail triangle $\{0, 2, k+2\}$ (entries $M_{04} = M_{24} = 1$),
and the second tail triangle $\{2, k+2, k+3\}$
(entries $M_{25} = M_{45} = 1$). Deleting $v_{\mathrm{tail}}$ removes
the last row and column, yielding a $5 \times 5$ matrix $M_{H}(m)$
for the deleted graph $H = \BTkt{k}{2} - v_{\mathrm{tail}}$.

By cofactor expansion along the $u_{P}$ row, the characteristic
polynomials are
\[
\det(\lambda I - M_{G}(m))
= \lambda^{6} - (2m + 7)\lambda^{4} - (2m + 6)\lambda^{3}
+ (7m + 3)\lambda^{2} + (10m + 2)\lambda + 3m,
\]
\[
\det(\lambda I - M_{H}(m))
= \lambda^{5} - (2m + 5)\lambda^{3} - (2m + 4)\lambda^{2}
+ 3m\lambda + 2m.
\]
The leading-order eigenvalues scale as $\pm\sqrt{2m}\,(1 + o(1))$
(the page-cluster modes coming from the $\sqrt{m}$ couplings to
$u_{P}$); the remaining interior eigenvalues converge to finite
limits.

The interior limit for $M_{G}$ is obtained by writing the eigenvector
as $(v_{0}, v_{1}, v_{2}, v_{P}, T_{1}, T_{2})$, using $v_{0} + v_{1} \to 0$
from the $u_{P}$-row at leading order, and substituting into the
remaining rows; the consistency reduces to the cubic
$P_{G}(\lambda) = 2\lambda^{3} - 7\lambda - 3 = 0$, plus a separate
limiting eigenvalue $\lambda \to -1$ as $m \to \infty$ from the
spine-antisymmetric direction (book-page swap; this is not an exact
eigenvalue of $M_{G}(m)$ at finite $m$). The interior limit for $M_{H}$ similarly reduces to
$P_{H}(\lambda) = 2\lambda^{3} + 2\lambda^{2} - 3\lambda - 2 = 0$.

By Vieta on $P_{G}$: $\sum \rho_{i} = 0$, $\sum \rho_{i} \rho_{j} = -7/2$,
so $\sum \rho_{i}^{2} = 7$. By Vieta on $P_{H}$:
$\sum \sigma_{i} = -1$, $\sum \sigma_{i}\sigma_{j} = -3/2$, so
$\sum \sigma_{i}^{2} = 4$. The trace identity gives
\[
\dminus(v_{\mathrm{tail}})
\;=\; 2 \deg_{\BTkt{k}{2}}(v_{\mathrm{tail}}) - \dplus(v_{\mathrm{tail}})
\;=\; 4 - \dplus(v_{\mathrm{tail}}).
\]
We argue that $\dplus(v_{\mathrm{tail}}) \to \alpha^{2} - \beta^{2}$,
where $\alpha$ (resp.\ $\beta$) is the unique positive root of
$P_{G}$ (resp.\ $P_{H}$).

The page cluster $B_{k}$ embedded in $\BTkt{k}{2}$ has its
spine-symmetric leading eigenvalues approaching $\pm\sqrt{2k}(1+o(1))$.
At the level of the squared difference
$\mu^{(G)\,2}_{\pm} - \mu^{(G-v)\,2}_{\pm}$, the naive cancellation
$\sqrt{2k} - \sqrt{2k} = 0$ is not enough: writing
$\mu^{(G)}_{\pm} = \pm\sqrt{2k} + a_{\pm} + o(1)$ and
$\mu^{(G-v)}_{\pm} = \pm\sqrt{2k} + b_{\pm} + o(1)$, the squared
difference equals $\pm 2\sqrt{2k}\,(a_{\pm} - b_{\pm}) + O(1)$, so we
need the next-order corrections to agree: $a_{\pm} = b_{\pm}$. This is
established by Cauchy interlacing applied to the rank-$1$ rooted
perturbation that distinguishes the two graphs --- $\BTkt{k}{2}$ vs.\
$\BTkt{k}{2} - v_{\mathrm{tail}}$ differ by the single ear $k+3$
attached via the edge $\{2, k+2\}$, and the leading book-spine
eigenvector has entries of order $1/\sqrt{k}$ at vertices $2$ and
$k+2$. By the secular equation for the bordering, the eigenvalue shift
is $\epsilon_{\pm} = \langle b, u_{\pm}\rangle\,/\,(\mu_{\pm} \mp \sqrt{2k}) \cdot O(1)
= O(1/k)$, hence $a_{\pm} - b_{\pm} = O(1/\sqrt{k})$ and the squared
difference of the leading eigenvalues is $o(1)$. The same control
extends to the equitable-partition quotient: $\BTkt{k}{2}$ has an
$S_{k-1}$-symmetry on its $k-1$ non-attached pages, giving a $6\times 6$
quotient matrix (cells $\{0\}, \{1\}, \{2\}, \{p_{3},\ldots,p_{k+1}\},
\{k+2\}, \{k+3\}$), and the deletion $\BTkt{k}{2} - v_{\mathrm{tail}}$
has a $5\times 5$ quotient. The two leading quotient eigenvalues of
each grow like $\pm\sqrt{2k} + O(1/\sqrt{k})$ with matching first
correction; their squared difference is $o(1)$. The page-antisymmetric
multiplicity-$(k-2)$ eigenvalue $0$ contributes $0$ to both $\splus$
and $\sminus$ and cancels exactly.

The remaining \emph{bounded} eigenvalues of the two quotients are the
content of the interior limit: rescaling the quotient matrices and
sending $k\to\infty$ produces an effective $4\times 4$ matrix
$M_{G}(\infty)$ for $\BTkt{k}{2}$ and $3\times 3$ matrix $M_{H}(\infty)$
for the deletion; their characteristic polynomials factor as
$(\lambda + 1) P_{G}(\lambda)$ and $P_{H}(\lambda)$, respectively.
Summing positive-root squares gives $\dplus(v_{\mathrm{tail}})
\to \alpha^{2} - \beta^{2}$, hence
$\lim_{k \to \infty} \dminus(v_{\mathrm{tail}}) = 4 - \alpha^{2} + \beta^{2}$.
Numerically $\alpha^{2} \approx 4.2294$, $\beta^{2} \approx 1.2647$, so
$4 - \alpha^{2} + \beta^{2} \approx 1.0353$.

The full derivation of the cubic resolvents and the matching
calculation is in \texttt{docs/lprime\_selector.md}; the result is
regression-locked in \texttt{tests/test\_lprime\_subfamilies.py}.

\emph{Status of this proof.} The page-cluster reduction, the cubic
resolvents $P_{G}, P_{H}$, the interior limits, and the numerical
limit value $4 - \alpha^{2} + \beta^{2} \approx 1.0353$ are derived
in full above and match the regression-locked numerical computation.
The leading-eigenvalue square-difference cancellation
($\pm \sqrt{2k}$ from the embedded $\Bk$ block) is presented above
modulo the secular-equation rate
$a_{\pm} - b_{\pm} = O(1/\sqrt{k})$, which is sketched at the level of
$\langle b, u_{\pm}\rangle/(\mu_{\pm} \mp \sqrt{2k}) \cdot O(1)$ but
not written in full detail in the paper; this is the conditional
subobligation under which the present proposition is stated. The
numerical limit is beyond doubt; the cleanly written cancellation
lemma is the only missing piece for promotion to an unconditional
theorem.
\end{proof}

\begin{corollary}[The universal ear-deletion lemma is false]\label{cor:universal-false}
The asymptotic value $4 - \alpha^{2} + \beta^{2} \approx 1.0353$ of
\Cref{thm:BT-asymptotic} is strictly below $17/16 = 1.0625$ by
$\approx 0.0272$. Hence the statement ``every simplicial degree-$2$
ear $v$ of every $2$-tree on $n \ge 4$ vertices satisfies
$\dminus(v) \ge 17/16$'' is false, and is false by a bounded constant
in the asymptotic regime.
\end{corollary}

\Cref{cor:universal-false} motivates the move to the \emph{existential}
form $\Lprime$ in \Cref{sec:lprime} and the \emph{max-degsum} selector
for picking the witness ear.

\subsection{$2$-paths at finite $n$: a Demmel--Kahan certificate}
\label{sec:subfamily:2path-finite}

\begin{theorem}[Engineering-rigorous floor on $\dminus(\Ln)$ for $n \le 1000$; conditional on the \texttt{dsyevd} backward-error envelope]
\label{thm:2path-DK}
Assume the conservative working envelope $C_{n} \le 10^{4}\, n$ for
the forward eigenvalue error of LAPACK's \texttt{dsyevd} symmetric
eigensolver, as discussed in Step~2 of the proof below
(\Cref{rem:fp-vs-interval}). Under this envelope, for every
$n \in [4, 1000]$,
\[
\dminus(\Ln) \;\ge\; \tfrac{17}{16} + \tfrac{1}{4} \;=\; \tfrac{21}{16}
\;>\; \tfrac{17}{16},
\]
with an empirical $50$-digit \texttt{mpmath} confirmation on a curated
$17$-record subset (\Cref{rem:fp-vs-interval}, \Cref{app:repro:DK}).
\end{theorem}

\begin{proof}
We give a constructive certificate.

\paragraph{Step 1 (floating-point computation).}
For each $n \in [4, 1000]$, compute the eigenvalues
$\tilde\lambda_{1}, \ldots, \tilde\lambda_{n}$ of the pentadiagonal
Toeplitz $\Adj(\Ln)$ in IEEE double precision via Householder
tridiagonalisation followed by symmetric divide-and-conquer (the
LAPACK routine \texttt{dsyevd}, which is the default driver behind
\texttt{numpy.linalg.eigvalsh}). Set
$\tilde s_{n}^{-} := \sum_{i:\, \tilde\lambda_{i} < 0} \tilde\lambda_{i}^{2}$
and $\tilde\delta_{n}^{-} := \tilde s_{n}^{-} - \tilde s_{n-1}^{-}$.

\paragraph{Step 2 (Demmel--Kahan-style a-posteriori bound, conservative envelope).}
For a backward-stable symmetric eigensolver the forward eigenvalue
error obeys an inequality of the form
$|\tilde\lambda_{i} - \lambda_{i}| \le C_{n}\, \epsilon_{\mathrm{mach}}\,
\|\Adj(\Ln)\|_{2}$, where $\epsilon_{\mathrm{mach}} \approx 2.22 \cdot 10^{-16}$
is the IEEE machine epsilon. The standard backward-error analyses of
Demmel \cite[Thm.~5.5]{DemmelANLA}, Parlett
\cite[\S 8.7]{Parlett1998}, and Anderson--Bai--Demmel
\cite[\S 4.7]{LAPACKUsersGuide} establish bounds of the form
$C_{n} \le c\, n$ for an absolute constant $c$ depending on the
specific algorithmic pipeline; published constants for symmetric
tridiagonal QR are of order $10^{2}$, and the
Householder-plus-divide-and-conquer
pipeline of LAPACK's \texttt{dsyevd} (the
default driver behind \texttt{numpy.linalg.eigvalsh}) plausibly
adds further low-order polynomial factors. We do \emph{not} attempt
in this paper a fully formal derivation of an explicit envelope
constant for the exact \texttt{dsyevd} path; the reference
implementation in \texttt{scripts/mpmath\_certify.py} uses
$c = 10$ as a safe upper bound (consistent with published constants
for backward-stable symmetric eigensolvers), and the certified
forward-error bounds below are stated under that script-default
envelope $c \le 10$. We note that the original analysis carried out
the same calculation under the deliberately loose working envelope
$c \le 10^{4}$; since any forward-error bound valid for $c \le 10^{4}$
is \emph{a fortiori} valid for $c \le 10$, the certificate holds
under either choice, and we exhibit the looser $c \le 10^{4}$ bound
in the displayed inequalities below for transparency and so that the
reader can see that the slack dwarfs the error even under the most
conservative envelope. Taking $C_{n} \le 10^{4}\, n$
and $\|\Adj(\Ln)\|_{2} \le \|f\|_{\infty} = 4$, the per-eigenvalue
forward error is at most
\[
|\tilde\lambda_{i} - \lambda_{i}| \;\le\; 10^{4}\, n \cdot 2.22 \cdot 10^{-16} \cdot 4
\;=\; 8.88 \cdot 10^{-12}\, n
\;\le\; 8.88 \cdot 10^{-9}
\qquad (n \le 1000).
\]
Propagating into $s_{n}^{-} = \sum_{i:\, \lambda_{i} < 0} \lambda_{i}^{2}$
via $|\delta(\lambda_{i}^{2})| \le 2 |\lambda_{i}|\, |\delta\lambda_{i}|
\le 2 \|\Adj(\Ln)\|_{2} \cdot 8.88 \cdot 10^{-12} n$ and summing over
at most $n$ negative eigenvalues,
\[
|\tilde s_{n}^{-} - s_{n}^{-}|
\;\le\; n \cdot 2 \cdot 4 \cdot 8.88 \cdot 10^{-12}\, n
\;=\; 7.1 \cdot 10^{-11}\, n^{2}
\;\le\; 7.1 \cdot 10^{-5}
\qquad (n \le 1000).
\]
Hence
\[
|\tilde\delta_{n}^{-} - \dminus(\Ln)|
\;\le\; |\tilde s_{n}^{-} - s_{n}^{-}|
+ |\tilde s_{n-1}^{-} - s_{n-1}^{-}|
\;\le\; 1.42 \cdot 10^{-4}.
\]

\paragraph{Step 3 (numerical floor + slack comparison).}
Direct numerical computation shows $\tilde\delta_{n}^{-} - 17/16 \ge 0.257$
for every $n \in [4, 1000]$, with the minimum $\approx 0.257$ attained
at $n = 6$ (the parity oscillation in $\tilde\delta_{n}^{-}$ has its
worst-case dip there). The slack $0.257$ dominates the certified
forward-error bound $1.42 \cdot 10^{-4}$ by more than three orders of
magnitude:
\[
\dminus(\Ln) \;\ge\; \tilde\delta_{n}^{-} - 1.42 \cdot 10^{-4}
\;\ge\; \tfrac{17}{16} + 0.257 - 1.42 \cdot 10^{-4}
\;>\; \tfrac{17}{16} + \tfrac{1}{4},
\]
as claimed.

The certificate is regression-locked in
\texttt{tests/test\_mpmath\_certify.py} and
\texttt{tests/test\_two\_path\_finite\_n.py}, with intermediate values
recorded in \texttt{data/two\_path\_mpmath\_certificate.json}.
\end{proof}

\begin{remark}[Failure mode F8: floating-point vs interval arithmetic]
\label{rem:fp-vs-interval}
The certificate in \Cref{thm:2path-DK} is engineering-rigorous rather
than interval-rigorous in two senses. First, the eigenvalues are IEEE
floating-point values, not certified intervals. Second, the constant
script-default envelope $c \le 10$ (and the looser $c \le 10^{4}$
working envelope used in the displayed analysis) in the
Demmel--Kahan-style bound is consistent with published constants for
backward-stable symmetric eigensolvers, but is not formally derived
for the exact LAPACK \texttt{dsyevd} pipeline within this paper; cf.\
\Cref{app:repro:DK}. The propagated error bound
$1.42 \cdot 10^{-4}$ on $\dminus(\Ln)$ is therefore the bound under
this conservative working envelope, not under a fully verified
constant. Crucially, the empirical slack $\ge 0.257$ exceeds this
propagated bound by more than three orders of magnitude, so any
plausible sharper constant (in particular the script-default
$c \le 10$, which gives a bound smaller by three orders of magnitude
than the displayed $c \le 10^{4}$ analysis) yields the same conclusion.
A confirmatory
\texttt{mpmath} run at $50$ decimal digits on a curated $17$-record
subset of $n$ (including the worst-case $n = 6$) is regression-locked
by the test suite (the suite spot-checks $9$ cases at $30$ dps with
$\sim 25$ digits of headroom; the full $50$-digit certificate is in
\texttt{data/two\_path\_mpmath\_certificate.json}), and reproduces
$\dminus(\Ln)$ to $\ge 13$ digits of agreement with the floating-point
values; no anomalous case is found. A full interval-arithmetic
upgrade (via \texttt{mpmath.iv} or a Krawczyk-style verified
eigensolver), or a formal derivation of an explicit envelope constant
for the \texttt{dsyevd} path, is mechanical but beyond the scope of
this paper; it is listed as open subobligation O5c.1 in the workstream
documentation.
\end{remark}

\subsection{Stieltjes-transform identification on $2$-paths}\label{sec:subfamily:2path-stieltjes}

\begin{theorem}[The $2$-path Stieltjes-transform limit]\label{thm:2path-stieltjes}
With the symbol $f(\theta) = 2\cos\theta + 2\cos 2\theta$ and with
$\Hgraph_{n} := \Ln - 1 \cong L_{n-1}$, the spectral measure of
$\Adj(\Hgraph_{n})$ tested against $\earw = e_{1} + e_{2}$ converges
weakly as $n \to \infty$ to an absolutely continuous measure
$\mu_{w, \infty}$ on $[-9/4, 4]$ whose density on the negative-spectrum
interval $(-9/4, 0)$ is
\[
\rho_{w}(\lambda) \;=\; \frac{1}{\pi}\, \sin\bigl(\theta_{2}(\lambda) - \theta_{1}(\lambda)\bigr),
\]
where $\theta_{1}(\lambda) < \theta_{2}(\lambda)$ are the two preimages
of $\lambda$ under $f$ in $(0, \pi)$. Writing $\nu^{-}_{\infty}$ for
the restriction of $\mu_{w, \infty}$ to $(-\infty, 0)$,
\[
\lim_{n \to \infty} \Wm(\Ln, \maxdeg) = \nu^{-}_{\infty}(\mathbb{R}),
\qquad
\lim_{n \to \infty} \Monem(\Ln, \maxdeg) = \int \lambda\, d\nu^{-}_{\infty},
\qquad
\lim_{n \to \infty} \Mtwom(\Ln, \maxdeg) = \int \lambda^{2}\, d\nu^{-}_{\infty}.
\]
\end{theorem}

\begin{proof}
The argument has four steps.

\paragraph{Step 1 (Stieltjes transform).}
For $z \in \mathbb{C} \setminus \spec(\Adj(\Hgraph_{n}))$, set
$S_{n}(z) := \earw^{\top}(zI - \Adj(\Hgraph_{n}))^{-1} \earw
= \int (z - \lambda)^{-1}\, d\mu_{w, n}(\lambda)$. The resolvent
matrix elements $G_{n}(z; i, j) := \langle e_{i}, (zI - \Adj(\Hgraph_{n}))^{-1} e_{j}\rangle$
satisfy the four-term recursion
\[
z\, G_{n}(z; i, j) - G_{n}(z; i-2, j) - G_{n}(z; i-1, j) - G_{n}(z; i+1, j) - G_{n}(z; i+2, j) = \delta_{ij}
\]
away from the boundary. Substituting $G(i) = \xi^{i}$ in the
homogeneous equation gives the self-reciprocal quartic
$\xi^{4} + \xi^{3} - z\xi^{2} + \xi + 1 = 0$, which under $u = \xi + 1/\xi$
reduces to the quadratic $u^{2} + u - (z + 2) = 0$ with roots
$u_{\pm}(z) = (-1 \pm \sqrt{4z + 9})/2$. For $\Im z > 0$, exactly two
of the four roots $\xi_{1}(z), \xi_{2}(z)$ lie strictly inside the
unit disk; on the spectrum
$\xi_{1} \to e^{-i\theta_{1}}, \xi_{2} \to e^{+i\theta_{2}}$ via
analytic continuation (the branch choice is verified at one base
point and propagated by continuity).

\paragraph{Step 2 (boundary collapse).}
The boundary recursions at $i = 1, 2$ together with the
bounded-at-$\infty$ ansatz $G_{\infty}(z; i, j) = A_{j}(z) \xi_{1}(z)^{i} + B_{j}(z) \xi_{2}(z)^{i}$
give a $2 \times 2$ linear system. Using the characteristic identity
$\xi^{4} + \xi^{3} + \xi - z\xi^{2} = -1$ (the quartic rearranged),
the four resolvent matrix elements $G_{\infty}(z; i, j)$ for
$i, j \in \{1, 2\}$ collapse to
\[
S_{\infty}(z) \;:=\; G_{\infty}(z; 1, 1) + 2\, G_{\infty}(z; 1, 2) + G_{\infty}(z; 2, 2)
\;=\; s(z)^{2} + s(z) - p(z),
\]
where $s(z) := \xi_{1}(z) + \xi_{2}(z)$ and
$p(z) := \xi_{1}(z) \xi_{2}(z)$. By strong-resolvent convergence of
finite-section banded Toeplitz operators \cite{BottcherSilbermann},
$S_{n}(z) \to S_{\infty}(z)$ uniformly on compact subsets of
$\mathbb{C} \setminus [-9/4, 4]$.

\paragraph{Step 3 (Stieltjes inversion).}
With our convention
$S(z) = \langle w, (zI - A)^{-1} w \rangle = \int (z - \lambda)^{-1}\, d\mu_{w}(\lambda)$,
a standard computation gives
$\lim_{\epsilon \to 0^{+}} \Im\, (\lambda + i\epsilon - \lambda')^{-1} = -\pi \delta(\lambda - \lambda')$
in distributional sense, so the Stieltjes inversion formula reads
$\rho_{w}(\lambda) = -\tfrac{1}{\pi}\, \lim_{\epsilon \to 0^{+}} \Im\, S_{\infty}(\lambda + i\epsilon)$
(note the minus sign). On $\lambda \in (-9/4, 0)$ the two interior roots of the quartic limit
to $\xi_{1} \to e^{-i\theta_{1}}, \xi_{2} \to e^{+i\theta_{2}}$ on the
unit circle (the Joukowski map $u = \xi + 1/\xi$ sending each
unit-circle root $\xi = e^{i\theta}$ to $u = 2\cos\theta$). By Vieta on
the quadratic $u^{2} + u - (z + 2) = 0$, the two $u$-values satisfy
$u_{-} + u_{+} = -1$, i.e.\ $2\cos\theta_{1} + 2\cos\theta_{2} = -1$,
hence
\[
\cos\theta_{1} + \cos\theta_{2} \;=\; -\tfrac{1}{2}.
\]
Substituting into $S_{\infty} = s^{2} + s - p$ and using this constraint
simplifies the cross-terms to
$\Im\, S_{\infty}(\lambda + i 0) = -\sin(\theta_{2} - \theta_{1})$,
which combined with the minus sign in the inversion gives
$\rho_{w}(\lambda) = \sin(\theta_{2} - \theta_{1})/\pi$, non-negative
since $0 < \theta_{2} - \theta_{1} < \pi$.

Sanity check on the unsigned moments: integrating $\rho_{w}$ over the
full spectrum $[-9/4, 4]$ recovers $(M_{0}, M_{1}, M_{2}) = (2, 2, 7)$,
matching the walk-moment identity \Cref{lem:walk-moments} (at the
$2$-path endpoint, $|T_{ab}(\Hgraph_{n})| = 1$ and $\degsum = 2 + 3 = 5$,
giving $M_{2} = 5 + 2 = 7$).

\paragraph{Step 4 (moment convergence via Portmanteau).}
By the Portmanteau theorem \cite{Billingsley1968}, weak convergence
$\mu_{w, n} \to \mu_{w, \infty}$ implies $\int \phi\, d\mu_{w, n} \to
\int \phi\, d\mu_{w, \infty}$ for every bounded measurable $\phi$
whose discontinuity set has $\mu_{w, \infty}$-measure zero. For
$\phi_{k}(\lambda) := \lambda^{k} \mathbf{1}[\lambda < 0]$ with
$k \in \{0, 1, 2\}$: the discontinuity at $\lambda = 0$ is
$\phi_{0}(0^{-}) = 1, \phi_{0}(0) = 0$ (jump), while for $k \ge 1$
both one-sided limits and the value at $0$ are zero (no
discontinuity). The $\mu_{w, \infty}$-measure of $\{0\}$ vanishes by the
standard boundary computation: with our sign convention
$S(z) = \int (z - \lambda)^{-1}\, d\mu(\lambda)$, an atom of mass $a$
at $\lambda = 0$ would contribute $a/(i\epsilon) = -i a/\epsilon$ to
$S(i\epsilon)$, while the absolutely continuous part remains bounded
as $\epsilon \downarrow 0$. Hence
\[
\mu_{w, \infty}(\{0\})
\;=\; -\lim_{\epsilon \downarrow 0}\, \epsilon \cdot
\Im\, S_{\infty}(i\epsilon)
\]
(no $1/\pi$ factor; that factor belongs to the density Stieltjes
inversion of Step~3, not the atom formula). The point $\lambda = 0$
lies on the limiting spectrum (in fact $0 = f(\pi/3) = f(\pi)$, so
$0$ is a critical value of the symbol), so atom non-existence cannot
be read off from analyticity at $z = 0$ alone; we must check the
boundary limit, and in particular identify the two interior unit-disk
roots $\xi_{1}(z), \xi_{2}(z)$ as $z \to 0$ along the upper imaginary
axis.

Compute $S_{\infty}(z)$ at $z = i\epsilon$ as $\epsilon \downarrow 0$.
The characteristic quartic
$\xi^{4} + \xi^{3} - z\xi^{2} + \xi + 1 = 0$ at $z = 0$ factors as
$(\xi+1)^{2}(\xi^{2} - \xi + 1) = 0$, with double root $\xi = -1$ and
the two unit-circle roots $e^{\pm i\pi/3}$. With
$u_{\pm}(z) = (-1 \pm \sqrt{4z + 9})/2$ (principal branch of
$\sqrt{\cdot}$), $u_{+}(0) = 1$ gives the quadratic $\xi^{2} - \xi + 1 = 0$
with roots $e^{\pm i\pi/3}$ on the unit circle, and $u_{-}(0) = -2$
gives $\xi^{2} + 2\xi + 1 = (\xi + 1)^{2} = 0$ with double root
$\xi = -1$ on the unit circle.

\emph{Branch selection.} The two \emph{interior} unit-disk roots are
not both on the $u_{+}$ branch: the $u_{+}$-branch quadratic at
$z = 0$ has both roots on the unit circle (modulus exactly $1$), so it
is not strictly interior. We track the interior selection by
continuous deformation along the upper imaginary axis. Numerical
continuation from $z = i$ (where the branch is unambiguously chosen
as in Step~3) down to $z = i\epsilon$ for $\epsilon \in
\{10^{-1}, 10^{-2}, 10^{-4}, 10^{-6}\}$ yields, for each $\epsilon$, two
interior roots converging respectively to
\[
\xi_{1}(0^{+}) \;=\; -1 \qquad \text{(from the $u_{-}$ branch)},
\qquad
\xi_{2}(0^{+}) \;=\; e^{-i\pi/3} \qquad \text{(from the $u_{+}$ branch)};
\]
at $\epsilon = 10^{-6}$ for example, the numerical roots are
$\xi_{1} \approx -0.999592 + 4.08 \cdot 10^{-4} i$ and
$\xi_{2} \approx 0.500000 - 0.866025 i$ (the inside-root selector
exercised by \texttt{tests/test\_half\_line\_stieltjes.py} via
\texttt{numerical\_density\_at}). The selection is the
generic one: of the two reciprocal pairs $(\xi, 1/\xi)$ supplied by
the quadratic in $u$, the unique inside-disk member of each pair is
the analytic continuation from $z = i$ down to $z = i\epsilon$.
Substituting,
\[
s(0^{+}) \;=\; \xi_{1} + \xi_{2} \;=\; -1 + e^{-i\pi/3}
\;=\; -\tfrac{1}{2} - \tfrac{\sqrt{3}}{2}\, i,
\qquad
p(0^{+}) \;=\; \xi_{1}\xi_{2} \;=\; -e^{-i\pi/3}
\;=\; -\tfrac{1}{2} + \tfrac{\sqrt{3}}{2}\, i,
\]
\[
s(0^{+})^{2} \;=\;
\bigl(-\tfrac{1}{2} - \tfrac{\sqrt{3}}{2}\,i\bigr)^{2}
\;=\; -\tfrac{1}{2} + \tfrac{\sqrt{3}}{2}\, i,
\]
\[
S_{\infty}(0^{+}) \;=\; s^{2} + s - p
\;=\; \bigl(-\tfrac{1}{2} + \tfrac{\sqrt{3}}{2}\, i\bigr)
   + \bigl(-\tfrac{1}{2} - \tfrac{\sqrt{3}}{2}\, i\bigr)
   - \bigl(-\tfrac{1}{2} + \tfrac{\sqrt{3}}{2}\, i\bigr)
\;=\; -\tfrac{1}{2} - \tfrac{\sqrt{3}}{2}\, i,
\]
a finite \emph{complex} number (not real). The numerical evaluation at
$z = 10^{-6} i$ gives
$S_{\infty}(10^{-6} i) \approx -0.499851 - 0.866582 i$, agreeing with
the closed form $-1/2 - i\sqrt{3}/2$ to $\sim 10^{-4}$ and converging
in $\epsilon$ to that value.

The maps $\xi_{1}(z), \xi_{2}(z)$ extend continuously along the
upper-imaginary approach $z = i\epsilon$, $\epsilon \downarrow 0$
(the discriminant $4z + 9$ is non-zero at $z = 0$; the unique
inside-disk Joukowski preimage of $u_{\pm}(z)$ is continuous away
from $u = \pm 2$, and we hit the branch point $u_{-}(0) = -2$ only at
$z = 0$ itself, where the degenerate double-root limit $\xi \to -1$ is
the stated one-sided limit, approached continuously from inside the
disk). Consequently $\Im\, S_{\infty}(i\epsilon)$ is bounded as
$\epsilon \downarrow 0$, with limit $-\sqrt{3}/2$. Therefore
\[
\mu_{w, \infty}(\{0\})
\;=\; -\lim_{\epsilon \downarrow 0} \epsilon \cdot
\Im\, S_{\infty}(i\epsilon)
\;=\; -\lim_{\epsilon \downarrow 0} \epsilon \cdot
\bigl(-\tfrac{\sqrt{3}}{2} + o(1)\bigr)
\;=\; 0,
\]
i.e.\ $\mu_{w, \infty}$ has no atom at $0$. The bounded one-sided
density limit
$\rho_{w}(0^{-}) = \sin(\pi - \pi/3)/\pi = \sqrt{3}/(2\pi)$ from
Step~3 -- recovered consistently from the same boundary calculation
via $\rho_{w} = -(1/\pi)\Im S_{\infty}(\lambda + i 0^{+})$ -- supplies
the absolute continuity needed for the Portmanteau condition. Hence
Portmanteau applies and gives the stated moment convergence.
\end{proof}

The full numerical identification — including the closed form
$\Iinf{L} = (2(310\pi^{2} - 837\sqrt{3}\,\pi + 2187))/(27\pi(20\pi -
27\sqrt{3}))$ — and its consequence for condition~(a) of the
joint-invariant ansatz on the $2$-path family appear in
\Cref{sec:moment-form:stieltjes}.
